\documentclass[12pt,a4paper,oneside]{report}             % Single-side
%\documentclass[11pt,a4paper,twoside,openright]{report}  % Duplex

%\PassOptionsToPackage{chapternumber=Huordinal}{magyar.ldf}
\usepackage{t1enc}
%\usepackage[latin2]{inputenc}
\usepackage{amsmath}
\usepackage{amssymb}
\usepackage{enumerate}
\usepackage[thmmarks]{ntheorem}
\usepackage{graphics}
\usepackage{epsfig}
\usepackage{listings}
\usepackage{color}
%\usepackage{fancyhdr}
\usepackage{lastpage}
\usepackage{anysize}
\usepackage[magyar]{babel}
\usepackage[utf8]{inputenc}
\usepackage{sectsty}
\usepackage{setspace}  % Ettol a tablazatok, abrak, labjegyzetek maradnak 1-es sorkozzel!
\usepackage[hang]{caption}
\usepackage{hyperref}
\usepackage{float} % ezzel lehet helybe rakni az abrat
\usepackage{listings}


%--------------------------------------------------------------------------------------
% Main variables
%--------------------------------------------------------------------------------------
\newcommand{\vikszerzo}{Bakró Nagy István}
\newcommand{\vikkonzulens}{Hartmann Péter \& Reichardt András}
\newcommand{\vikcim}{Poros plazma kísérletek támogatása multiprocesszoros környezetben}
\newcommand{\viktanszek}{Szélessávú Hírközlés és Villamosságtan Tanszék}
\newcommand{\vikdoktipus}{Diplomatervezés 1. - beszámoló}
\newcommand{\vikdepartmentr}{Szélessávú Hírközlés és Villamosságtan Tanszék}

%--------------------------------------------------------------------------------------
% Page layout setup
%--------------------------------------------------------------------------------------
% we need to redefine the pagestyle plain
% another possibility is to use the body of this command without \fancypagestyle
% and use \pagestyle{fancy} but in that case the special pages
% (like the ToC, the References, and the Chapter pages)remain in plane style

\pagestyle{plain}
%\setlength{\parindent}{0pt} % áttekinthet?bb, angol nyelv? dokumentumokban
% jellemz? \setlength{\parskip}{8pt plus 3pt minus 3pt} % áttekinthet?bb, angol nyelv? dokumentumokban jellemz?
\setlength{\parindent}{30pt} % magyar nyelv? dokumentumokban jellemz?
\setlength{\parskip}{0pt}    % magyar nyelv? dokumentumokban jellemz?

\marginsize{35mm}{25mm}{15mm}{15mm} % anysize package
\setcounter{secnumdepth}{0}
\sectionfont{\large\upshape\bfseries}
\setcounter{secnumdepth}{2}
\singlespacing
\frenchspacing

%--------------------------------------------------------------------------------------
%	Setup hyperref package
%--------------------------------------------------------------------------------------
\hypersetup{
    bookmarks=true,            % show bookmarks bar?
    unicode=false,             % non-Latin characters in Acrobat’s bookmarks
    pdftitle={\vikcim},        % title
    pdfauthor={\vikszerzo},    % author
    pdfsubject={\vikdoktipus}, % subject of the document
    pdfcreator={\vikszerzo},   % creator of the document
    pdfproducer={Producer},    % producer of the document
    pdfkeywords={keywords},    % list of keywords
    pdfnewwindow=true,         % links in new window
    colorlinks=true,           % false: boxed links; true: colored links
    linkcolor=black,           % color of internal links
    citecolor=black,           % color of links to bibliography
    filecolor=black,           % color of file links
    urlcolor=black             % color of external links
}

%--------------------------------------------------------------------------------------
% Set up listings
%--------------------------------------------------------------------------------------
\lstset{
	basicstyle=\scriptsize\ttfamily, % print whole listing small
	keywordstyle=\color{black}\bfseries\underbar, % underlined bold black keywords
	identifierstyle=, 					% nothing happens
	commentstyle=\color{white}, % white comments
	stringstyle=\scriptsize\sffamily, 			% typewriter type for strings
	showstringspaces=false,     % no special string spaces
	aboveskip=3pt,
	belowskip=3pt,
	columns=fixed,
	backgroundcolor=\color{lightgray},
} 		
\def\lstlistingname{lista}	

%--------------------------------------------------------------------------------------
%	Some new commands and declarations
%--------------------------------------------------------------------------------------
\newcommand{\code}[1]{{\upshape\ttfamily\scriptsize\indent #1}}

% define references
\newcommand{\figref}[1]{\ref{fig:#1}.}
\renewcommand{\eqref}[1]{(\ref{eq:#1})}
\newcommand{\listref}[1]{\ref{listing:#1}.}
\newcommand{\sectref}[1]{\ref{sect:#1}}
\newcommand{\tabref}[1]{\ref{tab:#1}.}


\newcommand{\ud}{\ \mathrm{d}}
\renewcommand{\vec}[1]{\overrightarrow{\mathrm{#1}}}

\DeclareMathOperator*{\argmax}{arg\,max}
%\DeclareMathOperator*[1]{\floor}{arg\,max}
\DeclareMathOperator{\sign}{sgn}
\DeclareMathOperator{\rot}{rot}
\definecolor{lightgray}{rgb}{0.95,0.95,0.95}

\author{\vikszerzo}
\title{\viktitle}
\includeonly{
	%guideline,%
	project,%
	titlepage,%
	%declaration,%
	abstract,%
	introduction,%
	chapter1,%
	chapter2,%
	chapter3,%
	chapter4,%
	chapter5,%
	chapter6,%
	chapter7,%
	aftermatch
	%acknowledgement,%
	%appendices,%
}
%--------------------------------------------------------------------------------------
%	Setup captions
%--------------------------------------------------------------------------------------
\captionsetup[figure]{
%labelsep=none,
%font={footnotesize,it},
%justification=justified,
width=.75\textwidth,
aboveskip=10pt}

\renewcommand{\captionlabelfont}{\small\bf}
\renewcommand{\captionfont}{\footnotesize\it}

%--------------------------------------------------------------------------------------
% Table of contents and the main text
%--------------------------------------------------------------------------------------
\begin{document}
\singlespacing
%\include{guideline}

\pagenumbering{arabic}
\onehalfspacing
%--------------------------------------------------------------------------------------
%	The title page
%--------------------------------------------------------------------------------------
\begin{titlepage}
\begin{center}
%\includegraphics[width=60mm,keepaspectratio]{figures/BMElogo.png}\\
\vspace{0.3cm}
\textbf{Budapesti Mûszaki és Gazdaságtudományi Egyetem}\\
\textmd{Villamosmérnöki és Informatikai Kar}\\
\textmd{\viktanszek}\\[5cm]

\vspace{0.4cm}
{\huge \bfseries \vikcim}\\[0.8cm]
\vspace{0.5cm}
\textsc{\Large \vikdoktipus}\\[4cm]

\begin{tabular}{cc}
 \makebox[7cm]{\emph{Készítette}} & \makebox[7cm]{\emph{Konzulens}} \\
 \makebox[7cm]{\vikszerzo} & \makebox[7cm]{\vikkonzulens}
\end{tabular}

\vfill
{\large \today}
\end{center}
\end{titlepage}
\tableofcontents\vfill
%\include{declaration}
%--------------------------------------------------------------------------------------
% Feladatkiiras (a tanszeken atveheto, kinyomtatott valtozat)
%--------------------------------------------------------------------------------------
\clearpage
\begin{center}
\large
\textbf{FELADATKIÍRÁS}\\
\end{center}

\noindent A GPU kártyák teljesítményének növekedése magas potenciált jelent a 
szimulátorok készítésében.
Ezzel egyetemben más programozási metodika is kívánkozik hozzá.\\

\noindent A hallgató feladatának a következõkre kell kiterjednie:
\begin{itemize}
\item Válasszon egy problémát szimulálásra és mutassa be a modelljét
\item Mutassa be a CUDA számítási platformját és a programozási metodikáját
\item Készítsen szimulációs programot MATLAB és CUDA platformon
\item Vizsgálja meg a kettő számítási teljesítményét és számítási pontosságát
\item Adjon iránymutatást a szimulátor további fejlesztési irányához!
\end{itemize}

\begin{flushleft}
\vspace*{1cm}
\textbf{Tanszéki konzulens}: Reichardt András, egy. tanársegéd
\end{flushleft}

\begin{flushleft}
\vspace*{1cm}
Budapest, 2013. szeptember 9.
\end{flushleft}

%----------------------------------------------------------------------------
% Abstract in hungarian
%----------------------------------------------------------------------------
\chapter*{Kivonat}\addcontentsline{toc}{chapter}{Kivonat}

hablaty

\vfill

%\include{introduction}
%----------------------------------------------------------------------------
\chapter{A Poisson egyenlet bemutatása}
%----------------------------------------------------------------------------
\section{A Poisson egyenlet}
%----------------------------------------------------------------------------

A Poisson egyenlet egy elliptikus parciális differenciál egyenlet.
Az elektrosztatika, a mechanika, hővezetés és a fizika sok problémája erre az egyenletre vezet.\\
Általános alakja, ahol $\varphi$ és $f$ valós vagy komplex értékű vektorfüggvények:
\begin{align}
\nabla^2 \varphi &= f\\
\left( \frac{\partial^2}{\partial x^2} + \frac{\partial^2}{\partial y^2} + \frac{\partial^2}{\partial z^2} \right)\varphi(x,y,z) &= f(x,y,z).
\end{align}
$f = 0$ esetén a Laplace egyenletre redukálódik.
Általánosan elmondható, hogy a megoldást $\varphi$ nem kívánjuk tudni a teljes téren, csupán
egy jól meghatározótt tartományon belül ($\Omega \in \mathbb{R}$).
A differenciálegyenlet egyértelmű megoldásához a következő peremfeltételek érvényre juttatására van szükség:
\begin{itemize}
\item Dirichlet: $\varphi(x,y,z) = f_0(x,y,z) \quad (x,y,z) \in \partial\Omega$
\item Neumann $\frac{\partial\varphi}{\partial n}(x,y,z) = g_0(x,y,z) \quad (x,y,z) \in \partial\Omega$
\item stb.
\end{itemize}
Az így összeállt probléma már megoldható.
Megoldani szeparálással és más analitikus módszerrel illetve a következő numerikus módszerekkel lehetséges:
\begin{itemize}
\item Véges elemek
\item Véges differenciák ( ha a megoldás feltehetően elég sima )
\end{itemize}
\newpage

%----------------------------------------------------------------------------
\section{Az elektrosztatikus probléma}
%----------------------------------------------------------------------------
A dolgozat során az elektrosztatikus téret leíró, Poisson (Laplace) egyenletre
redukálódó egyenlet megoldását keresem.
Ismeretes, hogy ilyenkor áram nem folyik és a gerjesztő menniységek állandóak, 
azaz $\vec{J} = 0$ és $\frac{\partial}{\partial t} \equiv 0$. Ezen feltételekkel a Maxwell-egyenletek a következő alakot veszik fel:
\begin{align}
rot \vec{E} &= 0\\
div \vec{D} &= \rho\\
D &= \varepsilon \vec{E}
\end{align}
Az örvénymentes villamos tér felírhatő skalárpotenciál gradienseként:
\begin{equation}
\vec{E} = - grad \varphi
\end{equation}
Ezzel a lépéssel jutunk az elektrosztatika Poisson egyenletére:
\begin{align}
-div(\varepsilon grad \varphi ) &= \rho
\end{align}
Ha homogén az anyag akkor:
\begin{align}
	\nabla^2\varphi  &= \frac{\rho}{\varepsilon}
\end{align}































%----------------------------------------------------------------------------
\chapter{A Poisson egyenlet diszkretizálása és megoldási módszerei}
%----------------------------------------------------------------------------

\section{A probléma definiálása}


%----------------------------------------------------------------------------
\section{A háló illesztése}
%----------------------------------------------------------------------------


%----------------------------------------------------------------------------
\section{Deriváltak közelítése és az egyenlet megoldása}
%----------------------------------------------------------------------------



%----------------------------------------------------------------------------
\chapter{Nvidia CUDA platform ismertetése}
%----------------------------------------------------------------------------

\section{CUDA-val szerelt GPU-k bemutatása}


\section{Architektúra és programozása}



\section{Mintaprogram: Két vektor összegzése blokkok és szálak használatával}

Az eszközön futó kód:
\begin{lstlisting}[frame=single]  % Start your code-block

// Az eszkozon futo mag.
__global__ void add(int *a, int *b, int *c) {
    int index = threadIdx.x + blockIdx.x * blockDim.x;
    c[index] = a[index] + b[index];
}
\end{lstlisting}
















%----------------------------------------------------------------------------
\chapter{Megoldók ismertetése és összehasonlítása}
%----------------------------------------------------------------------------

\section{Megoldó MATLAB környezetben}


\section{Megoldó CUDA-C nyelven}

\section{Összehasonlítás}


%\include{chapter6}
%\include{chapter7}
%----------------------------------------------------------------------------
\chapter{Összegzés}
%----------------------------------------------------------------------------

	Dolgozatomban bemutattam a porosplazma kísérletek apparátusát. A kísérlet során a kristályrácsba
	rendeződő részecskékről egy nagysebességű kamerával fényképek készülnek. A dolgozatomban ezen
	képeket, kellett feldolgoznom és a részecskék pozícióját detektálnom.
	A pozíciók a fizikai modell/szimuláció validálására szolgálnak.
	
	Ismertettem a részecske detektálásának módszerét szűrés és adaptív döntési küszöb használatával.
	Az elterjedt FIR Gauss szűrő helyett a hatékonyabb medián szűrőt javasoltam és alkalmaztam. A
	pozíció számítására a momentum módszert implementáltam, ami nagyobb számítási energiát igényel, de
	szubpixeles felbontást tudtam vele elérni. Konstatáltam, hogy az így kialakult program masszívan párhuzamosítható.
	
	Ezután áttekintettem az OpenCL keretrendszert, amit a párhuzamos program megírásának segítségére
	használtam. Az itt ismertetett megállapításokat figyelembe véve állítottam össze a párhuzamos
	program lépéseit, amit részleteztem is.
	
	Végül az elkészült programot CPU-n és GPU-n is futtatva a futási idejüket összevetettem és
	azonosítottam a gyorsulás forrását kitérve a processzormagra és a memóriájára.
	
	\section*{További feladatok:}
	\begin{itemize}
		\item A host-program real-time mérésbe helyezése egy producer-consumer sémájú szál megoldás
		alkalmazásával,
		\item Az eredmény grafikus  megjelenítése pl.: OpenGL használtatával,
		\item Az OpenCL szabvány által specifikált vektor műveletek támogatásának kiaknázása, ami az Intel
		Xeon PHI processzorkártyában rejlő teljesítményt ki tudná aknázni.
	\end{itemize}
	
	
%\include{acknowledgement}

\chapter{Függelék}


%----------------------------------------------------------------------------
\section{Matlab szimulátor}
%----------------------------------------------------------------------------
\texttt{feszko.m }
\begin{lstlisting}[frame=single]  % Start your code-block

clear all; close all; clc;

L = 1;

N = 500;
dx = L / N;

f0 = 0;
f1 = 1;

IT = 10000;

%%

fi = zeros(IT, N);
fi(:, 1) = f0 * ones(IT, 1);
fi(:, end) = f1 * ones(IT, 1);

T = zeros(N, N);
for n = 2:N-1
    T(n-1, n) =  0.5;
    T(n+1, n) =  0.5;
end
T(1,1) = 1;
T(end,end) = 1;

tic

for i = 2:IT
   fi(i, :) =  fi(i-1, :) * T;
end

toc

figure;
ax1 = subplot(211);
plot(linspace(0, L, N), fi(end, :)); grid;
xlabel('x');
ylabel('phi(x)');
title('Potencial eloszlas az iteracio vegen');

ax2 = subplot(212);
plot(linspace(0, L, N-1), diff(fi(end, :)) ./ dx);  grid;
xlabel('x');
ylabel('E(x)');
title('Elektromos ter az iteracio vegen');

linkaxes([ax1 ax2], 'x');

mfig(fi(1:1000:IT,1:4:N), 1:1000:IT,1:4:N)
\end{lstlisting}

\texttt{feszko2D.m }
\begin{lstlisting}[frame=single]  % Start your code-block

clear all; close all; clc;

X = 500;
Y = 500;

% mindig oszthato legyen 2vel
%global NX NY;
NX = 12;
NY = 12;

dx = X / NX;
dy = Y / NY;

f0 = 0;
f1 =1;


IT = 20000;

%%
ij2n = ((ones(NY,1) * (1:NX)) + ((0:NX:NX*(NY-1))' * ones(1,NX)))';


index = zeros(NX, NY);

index( 1: NX/4, NY/4+1:end) = -5;
index((3*NX/4)+1:end, 1: 3*NY/4) = -5;

index(1: 3*NX/4, 1) = -2;
index(3*NX/4, 1: 3*NY/4) = -6;
index(3*NX/4:end, (3*NY/4)+1) = -2;
index(end, (3*NY/4)+1 : end) = -6;
index((NX/4)+1:end, end) = -8;
index((NX/4)+1, NY/4:end) = -4;
index(1:NX/4, NY/4) = -8;
index(1, 1: NY/4) = -4;

index(3*NX/4,1) = -3;
index(3*NX/4,(3*NY/4)+1) = -3;

index((NX/4)+1,NY/4) = -7;
index((NX/4)+1,end) = -7;

index(1, 1:NY/4) = -10;
index(end, (3*NY/4)+1:end) = -11;


flip(index', 1)
flip(ij2n', 1)

%%
T = zeros(NX*NY, NX*NY);
[ii, jj] = find(index == 0);

for n = 1: length(ii)
    nn = ij2n(ii(n), jj(n));
    T(ij2n(ii(n)+1, jj(n)  ) , nn ) = 0.25;
    T(ij2n(ii(n)  , jj(n)+1) , nn ) = 0.25;
    T(ij2n(ii(n)-1, jj(n)  ) , nn ) = 0.25;
    T(ij2n(ii(n)  , jj(n)-1) , nn ) = 0.25;
end

valind = [ 1  1;
        0  1;
       -1  1;
        1  0;
        NaN NaN;
       -1  0;
        1 -1;
        0 -1;
       -1 -1];
        

% mi mivel
m = 1;
for i = 1:NX
    for j=1:NY
        if(-10 < index(i,j) & index(i,j) < 0 & index(i,j) ~= -5)
            mm(m, 1) = ij2n(i,j);
            mm(m, 2) = ij2n(i+valind(-index(i,j),1), j+valind(-index(i,j),2));
            m = m + 1;
        end
    end
end

f00 = find(reshape(index, 1, []) == -10);
f11 = find(reshape(index, 1, []) == -11);

%%
fi = zeros(IT, NX*NY);
% dirich
   for m= 1:length(f00)
       fi(1, f00(m)) = f0;
   end
   for m= 1:length(f11)
       fi(1, f11(m)) = f1;
   end

tic

for i = 2:IT
   fi(i, :) =  fi(i-1, :) * T;
   for m=1:length(mm)
       fi(i, mm(m, 1)) = fi(i, mm(m, 2));
   end
   % dirich
   for m= 1:length(f00)
       fi(i, f00(m)) = f0;
   end
   for m= 1:length(f11)
       fi(i, f11(m)) = f1;
   end
end

toc

fii = zeros(NX,NY, IT);
for i= 1:IT
    fii(:,:, i) = reshape(fi(i, :), NX, NY);
end



%%
fii( 1: NX/4, NY/4+1:end, end) = NaN;
fii((3*NX/4)+1:end, 1: 3*NY/4, end) = NaN;

figure;
ax1 = subplot(211);
contourf(fii(:,:,end)');
xlabel('x');
ylabel('y');
title('Potencial eloszlas az iteracio vegen');

ax2 = subplot(212);
contourf(diff(fii(:,:,end))');
xlabel('x');
ylabel('y');
title('Elektromos ter az iteracio vegen');

linkaxes([ax1 ax2], 'x');

%%


figm = figure('Renderer','zbuffer');
contourf(fii(:,:,2)'); title('0');
axis tight
set(gca,'NextPlot','replaceChildren');
% Preallocate the struct array for the struct returned by getframe
F(200) = struct('cdata',[],'colormap',[]);
% Record the movie
for j = 1:200
    fii( 1: NX/4, NY/4+1:end, j*100) = NaN;
    fii((3*NX/4)+1:end, 1: 3*NY/4, j*100) = NaN;
    surf(fii(:,:,j*100)');
    title(['it = ' num2str(j*100)]);
    xlabel('N'); ylabel('IT');
    
    %contourf(fii(:,:,j*100)'); 
    F(j) = getframe;
end


figure;
movie(F, 2)
\end{lstlisting}


\texttt{main.cu}
\begin{lstlisting}[frame=single]  % Start your code-block

#include <stdio.h>
#include <stdlib.h>
#include <time.h>
#include <string.h>

#include <helper_functions.h>
#include <helper_cuda.h>


extern "C"
void CPU(
    float *h_C,
    float *h_A,
    float *h_B,
    int vectorN,
    int elementN
);



#include "kernel.cuh"


float RandFloat(float low, float high)
{
    float t = (float)rand() / (float)RAND_MAX;
    return (1.0f - t) * low + t * high;
}



///////////////////////////////////////////////////////////////////////////////
// Data configuration
///////////////////////////////////////////////////////////////////////////////

const int VECTOR_N = 456;
const int ELEMENT_N = 4096;

const int    DATA_N = VECTOR_N * ELEMENT_N;

const int   DATA_SZ = DATA_N * sizeof(float);
const int RESULT_SZ = VECTOR_N  * sizeof(float);



///////////////////////////////////////////////////////////////////////////////
// Main program
///////////////////////////////////////////////////////////////////////////////
int main(int argc, char **argv)
{
    float *h_A, *h_B, *h_C_CPU, *h_C_GPU;
    float *d_A, *d_B, *d_C;
    //double delta, ref, sum_delta, sum_ref, L1norm;
    StopWatchInterface *hTimer = NULL;
    int i;


	getchar();

    printf("%s Starting...\n\n", argv[0]);

    // use command-line specified CUDA device, otherwise use device with highest Gflops/s
    findCudaDevice(argc, (const char **)argv);

    sdkCreateTimer(&hTimer);

    printf("Initializing data...\n");
    printf("...allocating CPU memory.\n");
    h_A     = (float *)malloc(DATA_SZ);
    h_B     = (float *)malloc(DATA_SZ);
    h_C_CPU = (float *)malloc(RESULT_SZ);
    h_C_GPU = (float *)malloc(RESULT_SZ);

    printf("...allocating GPU memory.\n");
    checkCudaErrors(cudaMalloc((void **)&d_A, DATA_SZ));
    checkCudaErrors(cudaMalloc((void **)&d_B, DATA_SZ));
    checkCudaErrors(cudaMalloc((void **)&d_C, RESULT_SZ));

    printf("...generating input data in CPU mem.\n");
    srand(123);

    //Generating input data on CPU
    for (i = 0; i < DATA_N; i++)
    {
        h_A[i] = RandFloat(0.0f, 1.0f);
        h_B[i] = RandFloat(0.0f, 1.0f);
    }

    printf("...copying input data to GPU mem.\n");
    //Copy options data to GPU memory for further processing
    checkCudaErrors(cudaMemcpy(d_A, h_A, DATA_SZ, cudaMemcpyHostToDevice));
    checkCudaErrors(cudaMemcpy(d_B, h_B, DATA_SZ, cudaMemcpyHostToDevice));
    printf("Data init done.\n");


    printf("Executing GPU kernel...\n");
    checkCudaErrors(cudaDeviceSynchronize());
    sdkResetTimer(&hTimer);
    sdkStartTimer(&hTimer);
    GPU<<<128, 256>>>(d_C, d_A, d_B, VECTOR_N, ELEMENT_N);
    getLastCudaError("scalarProdGPU() execution failed\n");
    checkCudaErrors(cudaDeviceSynchronize());
    sdkStopTimer(&hTimer);
    printf("GPU time: %f msecs.\n", sdkGetTimerValue(&hTimer));

    printf("Reading back GPU result...\n");
    //Read back GPU results to compare them to CPU results
    checkCudaErrors(cudaMemcpy(h_C_GPU, d_C, RESULT_SZ, cudaMemcpyDeviceToHost));


    printf("Checking GPU results...\n");
    printf("..running CPU scalar product calculation\n");

	sdkResetTimer(&hTimer);
    sdkStartTimer(&hTimer);
    CPU(h_C_CPU, h_A, h_B, VECTOR_N, ELEMENT_N);
	 sdkStopTimer(&hTimer);
    printf("cPU time: %f msecs.\n", sdkGetTimerValue(&hTimer));


    
    exit(EXIT_SUCCESS);
}
\end{lstlisting}


\texttt{kernel.cuh}
\begin{lstlisting}[frame=single]  % Start your code-block

#define IMUL(a, b) __mul24(a, b)



#define ACCUM_N 1024
__global__ void GPU(
    float *d_C,
    float *d_A,
    float *d_B,
    int vectorN,
    int elementN
)
{
    //Accumulators cache
    __shared__ float accumResult[ACCUM_N];

    for (int vec = blockIdx.x; vec < vectorN; vec += gridDim.x)
    {
        int vectorBase = IMUL(elementN, vec);
        int vectorEnd  = vectorBase + elementN;

        for (int iAccum = threadIdx.x; iAccum < ACCUM_N; iAccum += blockDim.x)
        {
            float sum = 0;

            for (int pos = vectorBase + iAccum; pos < vectorEnd; pos += ACCUM_N)
                sum += d_A[pos] * d_B[pos];

            accumResult[iAccum] = sum;
        }

        for (int stride = ACCUM_N / 2; stride > 0; stride >>= 1)
        {
            __syncthreads();

            for (int iAccum = threadIdx.x; iAccum < stride; iAccum += blockDim.x)
                accumResult[iAccum] += accumResult[stride + iAccum];
        }

        if (threadIdx.x == 0) d_C[vec] = accumResult[0];
    }
}
\end{lstlisting}

 

\listoffigures\addcontentsline{toc}{chapter}{Ábrák jegyzéke}
\listoftables\addcontentsline{toc}{chapter}{Táblázatok jegyzéke}

%\bibliography{mybib}
\addcontentsline{toc}{chapter}{Irodalomjegyzék}
\bibliographystyle{plain}



\label{page:last}
\end{document}
